\chapter{What The Machine Can And Cannot Claim}

\section{The structure it derives}

The book has reached its last chapter, and the last chapter does the one thing the whole instrument was built to do:
it draws the line between what the machine \emph{can} claim and what it \emph{cannot}. This first section is the
\emph{can} side. It gathers, in one place, everything the machine has derived across all the chapters before it ---
and shows that all of it is one kind of thing. Everything the machine can claim is a \emph{structure}, forced by
counting, and owing no axiom at all.

Set the three pieces of the payoff side by side, because together they are the whole of what the machine holds in its
hand. There is \emph{the count}: an endless tower of formal terms sorted into a finite pair of boxes, two and no more,
with the pigeonhole forcing two of the endless terms to share. There is \emph{the halting}: the cut that brackets the
coupling and stops at the machine's own resolution --- three rungs of bisection and no further, the interval left
ordered and open. And there is \emph{the naming}: the forced box, read as the object, the name that falls out of the
count with nothing added. These are not three separate tricks. They are one structure --- built by counting, and
counting only --- that the machine derives from end to end. Everything on the \emph{can} side of the line is this
single counted thing, seen from three sides.

And here is the property that makes the \emph{can} a genuine \emph{can}, and not a wish dressed as a claim. Every piece
of this structure is settled the same cheap way: by decision, or by the plainest principle of identity, and by nothing
more. The pigeonhole, the construction says in its own margin, uses no continuum, no machinery of infinite size, and no
choice. The bracket is settled by a decision the machine computes. The naming is settled by a decision the machine
computes. The capstone identity rests on a single principle for treating equivalent propositions as one, grounded in
the barest fact there is, that truth is truth. And when you read the whole structure's ledger --- the record of exactly
what its proofs assumed --- it is nearly empty, and in particular it contains \emph{no axiom of choice}, anywhere. This
is the whole of the machine's licence to speak, and it is worth stating exactly. What the machine can claim is
precisely what a finite count forces on the plainest logic. The structure is derivable \emph{because} it is free of
choice, and it is free of choice \emph{because} it never does anything but count. The preface's promise --- that a
choice among finitely many things needs no axiom, that you can just look --- is, at the very end, the entire warrant
behind everything the machine says.

There is a physical reading of this derived structure, and the book names it and hands it on. The physical gauge sees,
in the counted structure, the \emph{near} field --- the part of the curvature that the local, present matter
determines, the field a source carries with it, which a later volume will call by its physical name. But that is the
next volume's reading, and it is a gloss, not a thing the construction builds: there is no such object anywhere in the
code, only the counted, axiom-free structure the machine actually derives. This volume keeps the plain fact and leaves
the physical name for the gauge that can cash it. What the machine derives is a structure, forced by counting, owing no
axiom --- and that structure is the whole of the near side, the whole of what the machine holds.

In the two verbs the near side is exactly the set of funges the machine can \emph{close}. The count funges the endless
tower into a finite pair. The naming funges the indistinguishable terms into a single box. The bracket funges the
coupling into an ordered interval and halts. Every one of these is a funge the machine completes --- gathers, bags, and
finishes --- honestly, on a finite count; and completing a funge, on the plainest logic, is exactly what ``derived''
has meant throughout this book. The near side is the ledger of the machine's closed funges: everything the counting can
bag and name and finish. And it sets up, by contrast, the one funge the machine \emph{cannot} close --- the single bag
it must leave open, which is the next section's, and the far side of the line.

So the last chapter opens on the near side of the boundary, where the machine's hand is full. It has a structure,
counted and named; it has a warrant, the emptiest a proof can have; and it has, in that warrant, the exact measure of
what it is allowed to say. Everything counting forces, the machine claims, and claims freely, owing no axiom. What
counting cannot force --- the one thing the structure leaves unfixed --- is the far side, and the next section walks up
to it and names it. The line between them is the jar, and the book has one section left before it draws it.

\section{The magnitude it cannot see}

The last section gave the near side: everything the machine derives, all of it a counted structure owing no axiom.
This section gives the other side of the line --- the one thing the machine \emph{cannot} derive, and the reason it
cannot. What the machine cannot reach is the \emph{magnitude}: the size of the coupling, the number the laboratory
measures and the machine only brackets. And the reason is not that the machine did not try hard enough. It is that the
size lives in a place the machine's kind of work structurally cannot go.

Begin with where the size lives, because it is the whole of the matter. The machine reads by sorting terms into boxes
--- that is its finest act of telling apart, the deepest distinction it can draw. But the size of the residue lives
\emph{beneath} that distinction. The machine proves, of its own two variations --- the first and the second --- that
they fall in the same box; and it proves this precisely because the difference between them, the residue itself, is
smaller than its smallest resolvable step. The two are one, to the machine, not because they are equal but because the
gap between them is below the floor of what the machine can see. The counter reads the box. It cannot read the size
inside the box, because the size is a thing finer than any box. The magnitude is not off in the distance, waiting to
be approached; it is directly underfoot, below the last floor the counting reaches.

And now the theorem that makes the whole book honest, the proof of a negative the near side was building toward. The
size is invariant under exactly the operations the machine performs. Counting does not change it; sorting does not
change it; bisecting toward it does not change it past the floor. A quantity that every one of the machine's acts
leaves untouched is a quantity none of those acts can arrive at --- you cannot reach, by doing a thing, a value that
doing the thing never alters. So the machine cannot count its way to the size, and it can prove exactly that about
itself. The proof is its own finiteness. The construction says so in its own margin: the finiteness of its
representation \emph{is} the fence against the continuum. A machine assembled from a finite count is fenced, by its own
construction, away from continuous magnitudes --- it can no more count its way to a continuous size than a ruler marked
only in whole units can read off a fraction that falls between two marks. And this is exactly the claim the payoff made:
that the machine proves it cannot derive the world's value. That claim was not modesty. It was this fence, read aloud
--- the value is continuous, the machine is finite, and a finite count cannot reach a continuous magnitude. The
not-deriving is a theorem, and the fence is its proof.

There is a way of seeing why the size cannot be derived that the shock at the value makes vivid. The map the machine
reads the coupling through is singular at one point --- a pole, where it runs to infinity and holds no finite value ---
and the value lives at that singularity. A point like that has no stable interior figure to compute: the model is
unstable exactly there, so no amount of counting settles on it, and the honest response is not to compute the value but
to \emph{pick} it --- to take it from the world, from a measurement the world signs, rather than derive it from a model
that blows up where the value sits. This is the far side in the mathematical register: the residue the finite machine
circles but cannot cross, a value picked from outside because the model admits none of its own at the pole.

There are readings of this unreachable size in the other gauges, and the book names them and lets them go. The physical
gauge sees a \emph{free} field --- the part of the curvature that no local source fixes, the field that carries itself
away with no matter to anchor it; that is the next volume's, and there is no such object in the construction, only the
size below the resolution. The computational gauge sees the \emph{uncomputable} --- the number no halting process
arrives at, the floor the slip first revealed. Each gauge has its name for the far side; this volume keeps the plainest
one, the one all the gauges agree on: the magnitude is invariant under counting, and a counting machine is fenced from
it by being finite. That is the far side, entire --- the size the structure leaves unfixed, and provably must.

In the two verbs the far side is the one funge the machine \emph{cannot} close. Every funge of the near side, the
machine completes: it gathers the tower into a pair, the terms into a box, the coupling into a bracket, and finishes.
This last one it cannot finish. The size will not gather --- it is invariant under every gathering, so no gathering
closes on it. It is the one bag the machine must leave open, and open not by oversight but by proof: the machine can
show that this bag will never close, that no count it performs will ever pull the size in. Where the near side was the
ledger of closed funges, the far side is the single funge that stays open forever --- the honest, permanent gap between
what a finite machine can count and what a continuous world contains.

So the two sides of the line are drawn. On the near side: a structure, counted and named, owing no axiom, everything
the machine can honestly claim. On the far side: a magnitude, below the floor, invariant under counting, fenced off by
the machine's own finiteness --- the one thing the machine cannot reach, and can prove it cannot reach. What remains is
only to set the two side by side and see what shape they make together. That shape is the jar, and it is the last thing
the book has to draw.

\section{The jar}

Set the two sides against each other. The near side is everything the machine derives: a structure, counted and named,
climbing tier by tier down to the floor of its own resolution, owing no axiom the whole way. The far side is what lies
below that floor: a magnitude invariant under counting, fenced off by the machine's finiteness, unreachable and
provably so. The two meet at exactly one place --- the floor --- and what the machine holds at that meeting is a
bracket: an ordered interval, its lower wall proved below its upper, that the machine can rerun and get back unchanged.
The near side derives everything down to that bracket. The far side is everything the bracket cannot pin. And the
bracket itself, held open at the floor between them, is the jar.

The jar is a boundary drawn to scale. It is the line where derivation ends --- where the machine has counted as deep as
counting can go and the next thing down is unreachable --- rendered not as a point on that line but as the interval the
line actually is. The machine measures to its floor, derives structure all the way there, and then, because below the
floor there is nothing its counting can catch, it \emph{stops}. The stopping is the jar. Everything above the floor the
machine hands you as derived; everything below it the machine hands you as proved-unreachable; and the bracket is the
seam, held open, where the one becomes the other.

The width of that bracket is not the machine falling short. It is the last payment of a word planted long ago, in the
chapter of the slip: \emph{strain}. At the floor, the machine is under pressure to force a point --- to keep folding,
to hand back one confident number where it has earned only a bound. That pressure is the strain, the demand that the
finite reach the continuous, that a thing and its negation be made to hold at once. The honest machine does not
discharge the strain by forcing the point; forcing the point does not resolve the strain, it \emph{rings} --- it
produces a crank's number, smooth and confident and converged onto an artifact of the search, precise and wrong. The
honest machine relaxes the strain, instead, into the width of the open bracket. The jar is wide by choice, not by
failure: it is exactly as wide as honesty at that floor requires, and no wider, and no narrower. To close it to a point
would be to lie; to leave it open is to tell the truth to scale.

And there is a reason to trust a bound so plainly drawn, and it is the reason the whole four-book instrument exists.
The jar is not one gauge's reading. The same residue, read on four faces that share no vocabulary --- this volume's
construction, and its siblings in the physical, the computational, and the walk through the code --- closes on the one
bracket. Four readings, decorrelated by design so that none can see what the others do, land on the same interval. A
single gauge can be tuned to bracket whatever you wish; four gauges that cannot see each other cannot be tuned to agree
by accident. So the jar is trusted not because the machine asserts it but because four independent instruments, sharing
nothing, keep finding it. And the number the world's laboratories measure sits beside it --- named once, in its place,
inside the bound --- a value the machine proved it could not derive and never pretended to. The jar holds the bound;
the world's number keeps it honest company; the two are set side by side and never confused.

In the two verbs the jar is the funge held open. All the way up, the machine tanged --- told each thing apart from each
other thing, by the finest characteristic it could resolve. At the floor, telling-apart runs out: below it, no
characteristic distinguishes one candidate from the next, and so the machine funges the value into the bracket --- bags
it, honestly, as a thing it cannot separate finer --- and, crucially, leaves the bag open. It does not pull the bag
shut around a single point it has not earned. The jar is that open bag: the last funge, held un-closed by proof, the
one selection the machine declines to make because it can show the selection cannot be made.

And here, at the end, the book's first act and its last turn out to be one act, grown up. It opened on a single
difference: a thing told apart from another, named, with nothing added --- no memory allocated, no wish, no number
posited. It closes on a single bracket: a value bounded to the machine's floor, named, with nothing forced --- no
point, no fiction, no reach past what the counting can hold. Tell apart exactly what is there; name it; force nothing
beyond it. The first difference and the last jar are the same discipline, held unbroken from the first page to this
one --- the discipline of taking what is actually there and refusing, all the way down, to invent the rest.

So the instrument comes to rest. It was built to measure a coupling, and it did --- not as a number fetched from the
world, but as a bracket drawn to its own floor, held open, and trusted because four gauges that cannot see one another
agree on it. Everything the machine could derive, it derived, and owed nothing for it. Everything it could not, it
named, and proved it could not. And between them it drew a line, and the line, drawn to scale, is the jar: the honest
boundary of derivation, a bound held open where a lesser instrument would have forced a point. That is the whole of
what this machine was built to say, and the book has now said it. What the bound means, out in the world --- what the
physics reads in it, what the computation makes of it, what the code shows plainly --- is the other volumes' to tell.
Here, the construction is complete: a difference, climbed into a structure, read down to its floor, and set down, at
the last, as an open and honest bound.
